\section{Invited Presentations, Workshops, and Lectures}

Tatar, K. is the leading author of the following public presentations unless otherwise is indicated. \\

\cvitem{2025-10-22}{\textit{TRA 385 - Machine Learning and AI through Artistic Innovation}. I presented the pedagogical approach and the curriculm design of TRA 385 on TRACKs Teachers' Day in 2025.}

\cvitem{2025-05-16}{\textit{Tutorial on AI and ML for Artistic Innovation.} At the event, \textit{Redir - Spring meeting for principals and deans in science and technology (“Rektorer och dekaner i riket”)} at \href{https://www.chalmers.se/en/education/your-studies/course-selection-and-registration/select-courses/choose-a-tracks-course/}{TRACKs}, Chalmers.}

\cvitem{2024-03-20}{\textbf{Lecture.} \href{https://www.youtube.com/watch?v=HyRaBvoLvnI}{\textit{Organized Sound Spaces with Machine Learning}}, at \href{https://www.hfmt-hamburg.de/}{Hochschule für Musik und Theater Hamburg} in \textcolor{magenta}{Germany}, as a part of \href{https://www.youtube.com/playlist?list=PLlVs3aa4vQLpU4ni0sQExbGBT5YhKyWpg}{\_MUTOR Online Lecture Series - Artificial Intelligence for Music and Multimedia}.}

\cvitem{2024-03-20}{\textbf{Keynote.} \href{https://cinetic.arts.ro/en/evenimente/ai-in-art-practices-and-research-conference/}{AI in Art Practices and Research Conference.} \href{https://cinetic.arts.ro/en/about/}{CINETic - Centrul Internațional de Cercetare și Educație în Tehnologii Inovativ Creative},  I.L. Caragiale” National University of Theatre and Film in \textcolor{magenta}{Bucharest, Romania}.}

\cvitem{2023-11-15}{\textbf{Panel.} at WASP-HS - AI for Humanity and Society 2023, with Irina Shklovski, Guest Professor of Communication and Computing at Linköping University, Barry Brown Professor of Human-Computer Interaction at Stockholm University, Anne Kaun, Professor in Media and Communication Studies at Södertörn University. Malmo, Sweden}

\cvitem{2023-05-03}{\textbf{Seminar.} \href{https://www.communication.aau.dk/un-disciplinary-e73485#3rd--may}{\textit{Notes on Machine Learning and Artificial Intelligence for Artistic Practices}} as a part of \href{https://www.communication.aau.dk/un-disciplinary-e73485#3rd--may}{Un-Disciplinary} at Aalborg University in \textcolor{magenta}{Denmark}.}

\cvitem{2022-09-14 2022-09-17}{\textbf{Workshop.} \href{http://www.emutelab.org/blog/robot_opera}{\textit{Robots, AI, and Opera workshop}} at Aalborg University in \textcolor{magenta}{Denmark}.}

\cvitem{2022-09-13}{\textbf{Seminar.} \href{https://www.aau.dk/ai-and-creativity-e36349}{\textit{A Different Take on AI: Design Processes and Interaction in Machine Learning and Artificial Intelligence for Music and Arts}}, as a part of \textit{AI and Creativity: Exploring the Future of Music, Opera, and the Performing Arts in the Age of Machine Learning} at Aalborg University in \textcolor{magenta}{Denmark}.}

\cvitem{2022-06}{\textbf{Lecture}. \textit{An Introduction to Arts and Technology}. HCI Summer School 2022, organized by SIGCHI, Chalmers University of Technology, and Politechnika  Łódźka, hosted in \textcolor{magenta}{Łódź, Poland}}

\cvitem{2022-06}{\textbf{Workshop}.\textit{Sound Design using Deep Learning}. ACM HCI Summer School 2022, organized by SIGCHI, Chalmers University of Technology, and Politechnika  Łódźka, hosted in \textcolor{magenta}{Łódź, Poland}}

\cvitem{2022-03-31}{\textbf{Seminar and Panel}. \href{https://www.creative-ai-project.se/seminars/}{\textit{Musical Artificial Intelligence Architectures with Unsupervised Learning in Improvisation, Audio-Visual Performance, Interactive Arts, Dance, and Live Coding}} at the event \href{https://www.creative-ai-project.se/seminars/}{\textit{Interaction with generative music frameworks}}. As a part of seminar series \textit{dialogues: probing the future of creative technology} hosted virtually by \href{https://www.creative-ai-project.se}{Creative-Ai research group} of KTH Royal Institute of Technology, \textcolor{magenta}{Stockholm, Sweden}.}

\cvitem{2022-03-26}{\textbf{Panel}. Exploring links between body/image. K{\i}van\c{c} Tatar and Maria Pisiou. As a part of \href{https://4bidgallery.wordpress.com/2017/01/29/b-base/}{b.base expo} with Maria Pisiou at \href{https://4bidgallery.wordpress.com}{4bid Gallery} at OT301 in \textcolor{magenta}{Amsterdam, Netherlands}, invited through curatorial selection by 4bid gallery.}

\cvitem{2021-12}{\textbf{Seminar}. Incorporating Artificial Intelligence Architectures into Artistic Production and Live Performances, as a part of \href{https://www.sacmmt.com/}{Bowed Electrons 2021} hosted from \textcolor{magenta}{Cape Town, South Africa.}}

\cvitem{2021-10-15}{\textbf{Seminar}. Incorporating Artificial Intelligence Architectures into Artistic Production and Live Performances, at \href{https://kivanctatar.com/AI-Music-and-Improvisation-2021}{Pre-Conference - AI, Music and Improvisation} as a part of the conference Improvisation, Ecology and Digital Technology, hybrid format, in \textcolor{magenta}{Dusseldorf, Germany}}
\cvitem{2021-07-04}{\textbf{Artist Talk} on AI and Music Technology, with \href{https://remysiu.com/}{Remy Siu}, Yang Yi-Hsuan, and Su Li as a part of collective exhibition \href{https://event.culture.tw/NTMOFA/portal/Registration/C0103MAction?useLanguage=en&actId=10050&request_locale=en}{\textit{Aethereal}}, at the \href{https://www.ntmofa.gov.tw/en/}{Taiwan National Museum of Fine Arts}  \textcolor{magenta}{Taichung, Taiwan}}

\cvitem{2021-06-26}{\textbf{Workshop} on Creative Artificial Intelligence and Reinforcement Learning for Co-Creation, with \href{https://remysiu.com/}{Remy Siu}, as a part of collective exhibition \href{https://event.culture.tw/NTMOFA/portal/Registration/C0103MAction?useLanguage=en&actId=10050&request_locale=en}{\textit{Aethereal}}, at the \href{https://www.ntmofa.gov.tw/en/}{Taiwan National Museum of Fine Arts}  \textcolor{magenta}{Taichung, Taiwan}}
\cvitem{2020}{\textbf{Panel:} UNATC Distant Art CINETic Residencies. Tatar, K.; Tragtenberg, T.; Singer, M. O.; Hein, L. N.; and Schmitz C. \textcolor{magenta}{Ars Electronica 2020}, in \textcolor{magenta}{Linz, Austria}.}	

\cvitem{2020}{\textbf{Review Paper:} Review of Concert 12. \textit{Array: Special Issue - ICMC/NYCEMF 2019/ Reports+Reviews}. International Computer Music Association.}

\cvitem{2019-11}{\textbf{Workshop} on Creative Artificial Intelligence for Music and Multimedia, \href{https://cinetic.arts.ro/en/about/}{International Center for Research and Education in Innovative and Creative Technologies (CINETic), University of Theatre and Film ”I.L. Caragiale” (UNATC)}, \textcolor{magenta}{Bucharest, Romania}}

\cvitem{2019-02-08}{\textbf{Talk} at the Interactive Art, Science, and Technology in Western Canada (IAST Kelowna 2019) in conjunction with the \href{https://livingthingsfestival.com/}{Living Things 2019}, \textcolor{magenta}{Kelowna, BC, Canada}}

\cvitem{2018-06-01}{\textbf{Seminar}, \textit{Creative Artificial Intelligence for Music}, \href{http://www.miam.itu.edu.tr/}{Center for Advance Studies in Music (MIAM), Istanbul Technical University} \textcolor{magenta}{\.Istanbul, Turkey}}

\cvitem{2018-05-16}{\textbf{Seminar}, \textit{An Introduction to Metacreation and Musical Metacreation}, Music and Fine Arts Department, Middle East Technical University, \textcolor{magenta}{Ankara, Turkey}}

\cvitem{2018-02}{\textbf{Artist Talk} Modular approaches for Interactive Music Systems, at \textit{\href{https://www.facebook.com/events/2032925513589565}{Blue prints\_new prints}} at the Gold Saucer Studio in Vancouver, BC, Canada}

\cvitem{2017-09}{\textbf{Artist Talks} Three talks at the \href{https://www.aec.at/ai/en/iota}{Ars Electronica festival 2017 Artificial Intelligence}, \textcolor{magenta}{Linz, Austria}}

\cvitem{2017-06}{\textbf{Poster}, Tatar, K. \& Pasquier, P. (2017). MASOM: A Musical Agent Architecture based on Self-Organizing Maps, Affective Computing, and Variable Markov Models. Presented at the Eighth International Conference on Computational Creativity, ICCC 2017, \textcolor{magenta}{Atlanta, Georgia, USA}} 

\cvitem{2017}{\textbf{Demo}, \textit{Style Machine}, BC Tech Summit, Vancouver, BC, Canada}
\cvitem{2017}{\textbf{Demo}, \textit{Automatic Synthesizer Preset Generation with PresetGen}, BC Tech Summit, Vancouver, BC, Canada}

\cvitem{2016}{\textbf{Seminar}, \textit{An Introduction to Metacreation and Musical Metacreation}, Faculty of Computer and Informatics Engineering, Istanbul Technical University, \textcolor{magenta}{Istanbul, Turkey}}

\cvitem{2016}{\textbf{Workshop + Artist Talk}, \textit{Sound as a mediator in Interactive Arts: Challenges of Interdisciplinarity}, with Mirjana Prpa; at Vancouver ProMusica 2016, Vancouver, BC, Canada}

\cvitem{2016}{\textbf{Workshop + Artist Talk}, \textit{An interdisciplinary artwork: Pulse.Breath.Water}, with Mirjana Prpa; at \textit{movement.futures - May Residency 2016, School of Interactive Arts and Technology, Simon Fraser University, Vancouver, BC, Canada}}

\cvitem{2016}{\textbf{Workshop + Artist Talk}, \textit{The Potentials and Challenges  of VR, and Pulse.Breath.Water}, with Mirjana Prpa; at Chapel Sound Festival 2016, Vancouver, BC, Canada}

\cvitem{2016}{\textbf{Talk \& Demo} Prpa, M., Tatar, K., Pasquier, P., \& Riecke, B. E. (2016). \textit{Living In A Box: Potentials and Challenges of Existence in VR}, at \href{http://www.consumer-vr.com/}{the Consumer Virtual Reality (CVR) Conference, 2016}, Vancouver, BC, Canada}